\section{An intrinsic correction on conic-bundle nodes}\label{sf:section}

The surface calculation depends on a surface field inside the field of a
three-dimensional conic-bundle base.  We now remove that dependence.  All
fields and their inclusions in this section are over $\C$.  A function field
attached to a model is identified with its specified subfield of the common
function field; replacing the model does not replace that inclusion.
The test $u$ is the test of Section~\ref{pre:section}, for a collection of
Picard-number-one K3 surfaces of degree $d$ having no nonidentity
automorphism.  None of the bounds below depends on which subcollection is
used.  We will show that a nonzero correction singles out its embedded
surface field, and that this correction vanishes at the ends of links
over a point or curve.  These two assertions will make the surface
calculation telescope along a fourfold factorization.
We retain the convention of Section~\ref{sl:section} that ordinary
four-dimensional links belong to the paths chosen in
Theorem~\ref{bd:diagrams}\textup{(i)}.

\subsection{Eligible surface fields and the candidate}

Let $L$ be the function field of the three-dimensional base of a conic
Mori fibre space, and let $\alpha\in\operatorname{Br}(L)[2]$ be the
class of its generic conic.  The zero class is allowed.  A subfield
$k\subset L$ is \emph{eligible} if it is a surface function field and
$L$ is the function field of a smooth projective geometrically integral
genus-zero curve $C_k$ over $k$, with the following ramification bound.
Let $\Delta_k\subset C_k$ be the reduced finite subscheme of closed points
at which $\alpha$ has nonzero residue.  We require
\begin{equation}\label{sf:eligibility}
 r_k:=\deg_k\Delta_k\leq 7.
\end{equation}
Thus $r_k$ counts geometric points of the ramification support defined
\emph{before} extending the constants.  It does not mean the ramification
of the restriction of $\alpha$ after extension to an algebraic closure
of $k$.

The residue at each point of $\Delta_k$ defines a quadratic extension of
its residue field.  Together these give a degree-two finite \'etale cover
$\widetilde\Delta_k\to\Delta_k$.  If $r_k=3$, let $P_k$ be the finite
\'etale $k$-scheme whose geometric points choose one point in each of the
three fibres of this double cover.  It has degree eight.  Set
\begin{equation}\label{sf:candidate}
 \sigma_k(L,\alpha)
   =u(\Delta_k)+\mathbf{1}_{r_k=3}\,\frac12u(P_k).
\end{equation}
Here $u$ on a finite \'etale scheme is the sum of $u$ on its field factors,
as in the surface calculation.  The data in
\eqref{sf:candidate} are intrinsic to $(k\subset L,\alpha)$: a genus-zero
function field has a unique smooth projective curve model, and Brauer
residues are defined at its discrete valuations.

For the presentation supplied by a nontrivial generic two-ray move,
Proposition~\ref{sl:surface-identity} gives $r_k\leq7$.
Lemma~\ref{sl:conic-h} and Definition~\ref{sl:w-def} then identify
\eqref{sf:candidate} with the local correction $w-h_Q$.
The present definition also allows other eligible embedded fields;
we must compare their candidates before assigning a value to the node.

The half in \eqref{sf:candidate} causes no integrality ambiguity.
Complementing all three choices defines a fixed-point-free involution of
the geometric set $P_k$.  A field factor preserved by this involution
has a nonidentity automorphism.  If it were counted by $u$, this would
induce a nonidentity birational automorphism of its K3 minimal model;
birational maps between smooth minimal K3 surfaces are isomorphisms.
Such a factor is therefore not counted.  The counted factors occur in
pairs interchanged by complementation.  In particular $u(P_k)$ is even.
This observation is not needed for telescoping, but shows that all the
candidate corrections are nonnegative integers.

\subsection{Bounded ramification diagrams}

The following lemma supplies the implication needed in both comparison
arguments: a bounded diagram with marked ramification forces the
candidate to vanish at large degree.  We will then construct log paths
whose markings satisfy its hypothesis.  Boundedness includes the
morphisms and embedded marked support, rather than only the isomorphism
classes of their sources and targets, as in
Theorem~\ref{bd:diagrams}.

\begin{lemma}[Residues in a bounded diagram]\label{sf:residue-bound}
Suppose that a normal projective model of a genus-zero presentation
$L/k$ occurs in
a bounded family of diagrams
\[
                  Y\longrightarrow S,
\]
where $S$ is a normal projective surface model of $k$.  Suppose also
that the family has
bounded marked divisorial support containing the full visible
ramification of $\alpha$ on $Y$.  The family may instead be given on
the geometric generic fibre over a curve $Q$, in which case the
diagrams are over $Q$ and $\C(Q)\subset k$ is required.

For $r_k\leq7$, the orbit fields in $\Delta_k$, $\widetilde\Delta_k$,
and, when defined, $P_k$, are finite covers of degree at most fourteen
with branch contained in a uniformly bounded bad locus on the base
diagram.  Consequently none of these orbit fields is counted by $u$
once $d$ exceeds a constant depending only on the bounded family.
\end{lemma}

\begin{proof}
Resolve the diagram and marked support in bounded families, marking all
resolution exceptional divisors as well.  Write $D$ for the resulting
SNC support on the smooth total space.  It is enough that $D$ contain
the ramification; its unramified components do no harm.  Every prime
on the resolution is either the strict transform of a prime on the
old model or an exceptional prime.  Marking all exceptional primes
therefore includes any newly visible ramification, even over the
singular locus of the old model.

We form residues and their quadratic covers before extending constants:
the ground field $F$ is $\C$, or $\C(Q)$ on the ordinary generic fibre
over $Q$.  After finding their unramified open, we will base change
these same covers to $\overline{\C(Q)}$ to use the geometric bounds.
The Kummer
isomorphism identifies the Brauer $2$-torsion of the total
function field with its $H^2$ with coefficients in $\mu_2$
\cite[\S2, equations~(1)--(4)]{AuelBohningBothmerPirutka20}.  On the
smooth total model, the first two residue maps form a complex
\[
 H^2(F(Y),\mu_2)
 \xrightarrow{\partial}
 \bigoplus_{E\in Y^{(1)}}H^1(F(E),\Z/2)
 \xrightarrow{\partial}
 \bigoplus_{z\in Y^{(2)}}H^0(F(z),\mu_2^{-1}),
 \qquad \partial^2=0.
\]
Here $Y^{(j)}$ denotes the codimension-$j$ points.  These are the
coniveau differentials of
\cite[Example~2.1 and Proposition~3.9]{BlochOgus74}; the same complex
applies over the characteristic-zero function field of $Q$.
Consider a ramification component $D_i$ and a codimension-one point
$z$ of $D_i$ outside its intersections with the other components of
$D$.  Every term except that from $D_i$ is zero in the displayed
secondary residue at $z$.  Since $D_i$ is smooth, that term is its
ordinary DVR residue, so $\partial_z(\partial_{D_i}\alpha)=0$.
If the quadratic residue is represented by
$a\in F(D_i)^*/F(D_i)^{*2}$, this says that
$\operatorname{ord}_z(a)$ is even.  After multiplication by a square,
$a$ is a unit at the DVR, and adjoining its square root is unramified
there.  Normalize the indicated open of $D_i$ in the quadratic
algebra, componentwise in the split case.  This normalization is
finite, is unramified at every codimension-one point, and has regular
target.  Purity of the branch locus makes it \'etale throughout that
open \cite[Lemma~58.21.4, Tag~0BMB]{StacksPurityBranch2026}.

The horizontal ramification components are generically finite over
$S$.  Delete their nonfinite and branch loci on $S$, the images of their
pairwise intersections, the images of their intersections with the
remaining marked support, and the singular and resolution bad loci of
the diagram.  These images are proper closed subsets: over the generic
point of $S$ the distinct horizontal components are distinct points
of the smooth proper curve.  The deleted sets form a bounded embedded
boundary after stratifying the family.  Over the resulting smooth open
$U\subset S$, the horizontal union is finite \'etale, its components are
smooth and disjoint from the other marked divisors, and its residue
double cover is finite \'etale.  Their degrees over $U$ are $r_k$ and
$2r_k$.  The scheme of sections defining $P_k$ is then a finite \'etale
scheme of degree eight over $U$.  Thus all the fields in the statement
have the asserted degree and branch bounds.

These degree and branch bounds exclude the large-degree K3 fields.
First consider a bounded diagram over $\C$.
Fix an orbit field $E/k$ in the statement, and let $S_E\to S$ be
the normalization in $E$.  Take a general pencil in a uniformly bounded
very ample system on the resolved base surface $S$.  Let $h$ bound the
genus of its normalized general member and let $b$ bound its intersections
with the deleted boundary.  Each normalized component of the pullback
of that member to $S_E$ has degree $n\leq14$ and, by Hurwitz
\cite[Lemmas~53.12.2 and~53.12.4]{StacksRiemannHurwitz2026},
\begin{equation}\label{sf:hurwitz}
                 2g-2\leq n(2h-2)+b(n-1).
\end{equation}
If $u(E)=1$, a smooth projective model of $S_E$ is birational to a
test K3 surface.  Resolve the pulled-back pencil on this model and take
its Stein factorization.  Its connected general fibre is one of these
normalized components.  This contradicts the bound
$2g-2\geq\sqrt d$ in Lemma~\ref{bd:pencil} for large $d$.

When the family is given only on the geometric generic fibre of
$Q$, the same construction is performed over
$\overline{\C(Q)}$ on the base curve there.  Finite \'etale covers remain
finite \'etale after this extension, although they may split.  The same
degree and Hurwitz bounds therefore hold on every geometric component.
We retain the base changes of the original ramification support and
residue covers; we do not recompute the support from the extended
Brauer class, whose residues may have become zero.
For any orbit field $E/k$, the inclusions $\C(Q)\subset k\subset E$
give a rational map from its surface model to $Q$.
If $u(E)=1$, resolve that map on the corresponding K3 surface and take
its Stein factorization.  Then
the geometric general fibres of that map are precisely the components
whose genera were bounded.  Lemma~\ref{bd:pencil} again gives the
contradiction.  No assertion about the complexity of an unmarked
Brauer class is used.
\end{proof}

\subsection{The coefficient and the full visible boundary}

We use the single rational coefficient supplied by
Corollary~\ref{bd:choice-order}.  If $c_*>0$ denotes the uniform upper
bound obtained there from the ordinary conic-base diagrams, this choice is
\begin{equation}\label{sf:c-choice}
                  0<c<\min\{1/4,c_*\}.
\end{equation}
With $c$ fixed, take the log three-dimensional diagram bounds of
Theorem~\ref{bd:diagrams} and the terminalization bounds of
Theorem~\ref{bd:terminalization}.  Only then fix the common lower bound
on $d$ required by Lemma~\ref{sf:residue-bound} and
Theorem~\ref{bd:exclusions} for these families.  These choices are
uniform in the eligible fields, test subcollections, factorizations,
and their lengths.

Here are the boundary conventions needed to apply those results.
Given finitely many presentations of $L$, take a common smooth
projective model $Y$ resolving the rational maps to projective models
of their bases.  Resolve the ramification support and let $D_Y$ be
\emph{exactly} the reduced divisor of ramification visible on this
model.  Use the pair
\begin{equation}\label{sf:start}
                         (Y,\Phi),\qquad \Phi=cD_Y.
\end{equation}
The divisor $D_Y$ has SNC support.  All models and maps in the log
geography are taken with the transformed support; they do not extract
prime valuations from $Y$.  Their transformed support is therefore
the full ramification visible on them.  If a higher common start is
needed, use its full visible ramification again, including the
ramified exceptional divisors that have appeared.

This last convention is compatible with the discrepancy requirement
in Theorem~\ref{bd:diagrams}.  For the first blowup of a smooth centre
of codimension $\ell\geq2$ contained in $r\leq\ell$ boundary components,
the log discrepancy is $a(F,Y,cD_Y)=\ell-cr$.  Its ordinary discrepancy
is therefore $a(F,Y,cD_Y)-1=\ell-1-cr\geq1-2c>0$.
For an arbitrary exceptional prime, the SNC inequality
\eqref{bd:snc-discrepancy} gives $a(F,Y,cD_Y)\geq2(1-c)>1$.
On a higher smooth start $\pi:Y'\to Y$, the coefficient of $F$ in
$K_{Y'}+cD_{Y'}-\pi^*(K_Y+cD_Y)$ is
$a(F,Y,cD_Y)-1+c\epsilon_F$, where $\epsilon_F$ is one if $F$ is
ramified and zero otherwise.  This coefficient is positive in either
case.  A negative log MMP increases discrepancies for contracted primes.
Lemma~\ref{bd:higher-start} therefore recovers the same endpoints, and
the common-start clauses of Theorem~\ref{bd:diagrams} apply with the
fixed coefficient in \eqref{sf:c-choice}.

For an eligible $k$, the degree on the generic curve is
\begin{equation}\label{sf:negative-degree}
             \deg(K+\Phi)=-2+cr_k\leq-2+7c<-\tfrac14<0.
\end{equation}
Run the log MMP over a surface model of $k$.  The relative programme
ends in a Mori fibre space over a two-dimensional base birational to
that surface; its embedded base field is still $k$.  Indeed the
generic relative dimension is one, and the final generic contraction
is a genus-zero curve to its ground field.  Geometric integrality
makes this ground field $k$ rather than a proper finite extension.
These relative negative steps are also permissible steps in the
common-start log geography over projective bases.

\subsection{Uniqueness of the eligible field}

\begin{proposition}[Uniqueness from a nonzero candidate]
\label{sf:unique}
For $d$ above the uniform bound fixed after \eqref{sf:c-choice}, if
an eligible $k\subset L$ has
$\sigma_k(L,\alpha)\ne0$, then every eligible surface subfield of
$(L,\alpha)$ equals $k$ as an embedded subfield of $L$.
\end{proposition}

\begin{proof}
Let $k'\subset L$ be any other eligible field; its candidate need not
be nonzero.  Resolve the two presentations on a common start
\eqref{sf:start}.  By \eqref{sf:negative-degree}, its log MMP has an
output with base field $k$ and another with base field $k'$.  Connect
these outputs by the log Sarkisov links of
Theorem~\ref{bd:diagrams}, and write their successive nodes as
$Z_i/S_i$ for $0\leq i\leq N$.  The birational maps from the common
start identify every $\C(Z_i)$ with $L$; set $K_i=\C(S_i)\subset L$.
Thus $K_0=k$ and $K_N=k'$.  The class $\alpha\in\operatorname{Br}(L)[2]$
is fixed throughout this path, and the support of each boundary is its
full visible ramification on the corresponding model.

Every Mori base in this three-dimensional programme has dimension at
most two.  If a common base of a link has dimension two, the endpoint
bases also have dimension two and their contractions to the common
base are birational: they are connected-fibre contractions between
normal varieties of the same dimension.  The induced identifications
are the ones coming from $L$.  Such a link therefore preserves the
embedded surface field.

If $k'\ne k$, let $m$ be the first index for which $S_{m+1}$ is not
a surface with $K_{m+1}=k$.  Its common base $Q$ has dimension at most
one, by the preceding paragraph.  At each preceding node the data
$(k\subset L,\alpha)$ are unchanged.  Hence its proper generic curve,
residue covers, and candidate are the original $C_k$, the original
covers, and $\sigma_k(L,\alpha)$.  The relevant part of the path is
\[
\begin{tikzcd}[column sep=large,row sep=large]
 Z_m \arrow[rr,dashed,"\text{first departure}"] \arrow[d]
   && Z_{m+1}\arrow[d] \\
 S_m\arrow[dr] && S_{m+1}\arrow[dl] \\
 & Q
\end{tikzcd}
\qquad
\begin{gathered}
 \C(S_0)=\cdots=\C(S_m)=k\subset L,\\
 \dim Q\leq1.
\end{gathered}
\]
By Theorem~\ref{bd:diagrams}, on the
geometric generic fibre of $Q$ its whole marked diagram is uniformly
bounded, including the projection to the endpoint surface base.
All ramification points on the proper generic curve are visible as
divisors on this total model: their closures dominate the surface
base and have codimension one.  They are among the marked transforms
by the full visible convention.  If $Q$ is a curve, its function
field is contained in $k$ because the endpoint base contracts to
$Q$.  Lemma~\ref{sf:residue-bound} applies and gives
$\sigma_k(L,\alpha)=0$, a contradiction.  Thus $k'=k$.
\end{proof}

\subsection{Vanishing at ends over smaller common bases}

At ends of ordinary four-dimensional links over a point or curve,
we must exclude both the del Pezzo correction on a surface base and
the eligible-field candidates on a three-dimensional conic base.
The surface-base case uses the bounded finite sets already defined;
the conic case will require a second log MMP output.

\begin{lemma}[Vanishing at surface bases]\label{sf:surface-end}
At an end $V\to B$ with $\dim B=2$ of an ordinary four-dimensional
link over a common base $Q$ of dimension at most one, the
degree-five or degree-six del Pezzo correction $w$ is zero for
uniformly large $d$.
\end{lemma}

\begin{proof}
The relative generic diagram over $Q$ is bounded by
Theorem~\ref{bd:diagrams}.  On the smooth del Pezzo locus, let
$a\in\{5,6\}$ be the fibre degree.  The vanishings
$H^1(\mathcal O)=H^2(\mathcal O)=0$ make the relative Picard scheme
\'etale
\cite[Theorem~4.8, Corollary~5.13 and Proposition~5.19]{KleimanPicard057}.
We identify its finite part that parametrizes conic-pencil classes.
For a line-bundle class $F$ on a geometric fibre $S$, Riemann--Roch gives
\[
 \chi\bigl(\mathcal O_S(F-mK_S)\bigr)
 =1+\tfrac12\bigl(F^2-(2m+1)F\cdot K_S+a m(m+1)\bigr).
\]
Consequently the equations $F^2=0$ and $-K_S\cdot F=2$ are equivalent
to having the fixed anticanonical Hilbert polynomial
$\tfrac a2m(m+1)+2m+2$.
This fixed-polynomial locus in the relative Picard scheme is open and
closed and of finite type; in the smooth projective family it is also
projective
\cite[Theorem~6.20, Exercise~5.7 and its answer]{KleimanPicard057}.
It is therefore proper and quasi-finite, hence finite \'etale.
Lemma~\ref{sl:pencil-sets} identifies its geometric fibres
with $A(S)$, of cardinality five or three according as $a=5$ or $6$.
These are statements about the Picard scheme and require no universal
line bundle on the original family.
The complement of this smooth family is bounded as an embedded locus
of the diagram.  Each field factor of the finite scheme is
therefore a cover of bounded degree unramified outside this bounded
locus.  Apply the bounded-cover exclusion of
Theorem~\ref{bd:exclusions}, or the Hurwitz calculation
\eqref{sf:hurwitz}.  When $Q$ is a curve its field is contained in
each tested surface field.  Thus no factor is counted by $u$ and
$w=0$.
\end{proof}

Now let $V\to T$ be a conic end with $\dim T=3$ and
common base $Q$ of dimension at most one, and put $F=\C(Q)$.
Corollary~\ref{bd:choice-order} gives a bounded family of klt log-Fano
pairs $(T_{\overline F},c(D_T)_{\overline F})$ with the retained marked
support and the fixed coefficient \eqref{sf:c-choice}.  These are the
inputs to the application of Theorem~\ref{bd:terminalization} in the
following proof.  The corollary proves
rank one over $F$; it requires no rank-one assertion over $\overline F$.

\begin{proposition}[Vanishing at conic ends]\label{sf:conic-end}
Let $V\to T$ be a conic end of an ordinary four-dimensional Sarkisov
link with common base $Q$ of dimension at most one.  For the uniform
large degrees under consideration, no eligible surface subfield of
$(\C(T),\alpha)$ has nonzero candidate.
\end{proposition}

\begin{proof}
Put $L=\C(T)$ and suppose an eligible $k\subset L$ has nonzero
candidate.  The bounded diagram on $T$ does not yet include the
presentation over $k$; in particular, $k$ need not contain $\C(Q)$.
We will construct an output with bounded marked geometric generic
diagram over $Q$ and compare it with the output over $k$.
Choose a common smooth
SNC start $Y$ mapping to $T$ and to a projective surface model of
$k$, with the full visible boundary \eqref{sf:start}.  One log MMP
output has base field $k$ by \eqref{sf:negative-degree}.

\paragraph{A bounded output over $Q$.}
First run the log MMP \emph{over $T$} to a relative minimal model
$b:Y^*\to T$.  This is the birational case: relative bigness holds,
so \cite[Theorem~1.2]{BCHM10} applies.  Let $D_*$ be the transform
of $D_Y$ on $Y^*$; its pushforward under $b$ is $D_T$.

We first study the geometric generic fibre over $Q$, retaining the
same symbols for these restrictions.  The marked supports
are the base changes of the original supports, regardless of whether
the extended Brauer class still ramifies on them.  We have
\begin{equation}\label{sf:exceptional-nef}
 K_{Y^*}+cD_* = b^*(K_T+cD_T)+E.
\end{equation}
The divisor $E$ is $b$-exceptional and $b$-nef; the latter follows
from relative nefness of the left side.  The negativity lemma gives
$E\leq0$.  For an exceptional prime $F$ of $b$, its coefficient is
\begin{equation}\label{sf:cutoff}
 \operatorname{coeff}_F(E)
 = a(F,T,cD_T)-1+c\epsilon_F,
 \qquad
 \epsilon_F=\begin{cases}1&F\text{ occurs in }D_*,\\0&F\text{ does not occur in }D_*.
 \end{cases}
\end{equation}
Consequently every valuation extracted by $b$ on this geometric
generic fibre satisfies
$a(F,T,cD_T)\leq1$; an extracted valuation occurring in $D_*$ even satisfies
$a(F,T,cD_T)\leq1-c$.

Theorem~\ref{bd:terminalization} now bounds the geometric generic
model and all its subsequent non-extracting contraction diagrams,
including the marked support.  Indeed the only extracted divisors
are among the bounded terminalization data for the bounded klt log
Fano pairs $(T,cD_T)$.  On the geometric generic fibre over $Q$,
every component of $D_*$ is either the strict transform of a component
of $D_T$, or one of these extracted prime divisors.  The markings here
are retained under extension of constants.  The terminalization
theorem marks the entire extraction list, including primes with
crepant coefficient zero, so it bounds this full support and all of
its subsequent transforms on that geometric generic fibre.  Primes
of the resolution exceptional over $T$ there and absent from this
list cannot survive on $Y^*$ there by \eqref{sf:cutoff}.
This bounds the geometric generic model $Y^*$ over $Q$ and its
subsequent non-extracting diagrams, not the arbitrary resolution $Y$.
To obtain a Mori output over $Q$, we verify
that the log canonical divisor on this geometric generic fibre is not
pseudo-effective.  To see this from \eqref{sf:exceptional-nef}, choose
a general complete-intersection curve in $T$ avoiding the centres
of the exceptional divisors of $b$.  Its strict transform has
negative intersection with $K_{Y^*}+cD_*$, because $K_T+cD_T$ is
anti-ample.  Such curves form a movable covering family, so a
pseudo-effective divisor could not have this negative intersection.

Return to the models over $\C$.  We may therefore run the log MMP
from $Y^*$ over $Q$ to a Mori fibre space $Z/B$.  Its geometric
generic diagram over $Q$ is bounded by the terminalization theorem,
since these additional steps do not extract divisors.

\paragraph{Comparison with the output over $k$.}
Connect the output with field $k$ to $Z/B$ by log links from the same
initial pair.  These links are taken in the absolute common-start
geography; the arbitrary eligible field $k$ need not contain $\C(Q)$.
If $B$ is not a surface with embedded field $k$, consider the first
departure from $k$ along this path, as in Proposition~\ref{sf:unique}.
Call the common base of that link $Q'$.  It has dimension at most one,
its incoming endpoint has the original candidate $\sigma_k(L,\alpha)$,
and, when $Q'$ is a curve, its morphism to $Q'$ gives
$\C(Q')\subset k$.  The bounded diagram over $Q'$ therefore forces
that candidate to vanish by Lemma~\ref{sf:residue-bound}.

If instead $B$ is a surface with $\C(B)=k$, the bounded geometric
generic diagram of $Z/B$
over the original $Q$ gives the same contradiction directly.  Here
$B\to Q$ gives $\C(Q)\subset k$ when $Q$ is a curve, so the other
application of Lemma~\ref{sf:residue-bound} has its required field
inclusion as well.  Both alternatives contradict the nonzero candidate.
\end{proof}

\subsection{The intrinsic function and telescoping}

We now attach the correction to a four-dimensional Mori node $V\to B$.
When $\dim B=3$, let
$\alpha\in\operatorname{Br}(\C(B))[2]$ be the class of the generic
conic.  Define a function on the entire node by
\begin{equation}\label{sf:node-function}
 s(V/B)=
 \begin{cases}
  0,&\dim B\leq1,\\
  w(V/B),&\dim B=2,\\
  \sigma_k(\C(B),\alpha),&\dim B=3\text{ and some eligible }k
       \text{ has nonzero candidate},\\
  0,&\dim B=3\text{ and no such }k\text{ exists}.
 \end{cases}
\end{equation}
The function $w$ in base dimension two is the del Pezzo function
of Definition~\ref{sl:w-def}.

\begin{theorem}[Intrinsic correction]\label{sf:intrinsic}
For uniformly sufficiently large $d$, the function
\eqref{sf:node-function} is well defined.  At a conic node its value
equals the candidate of \emph{every} eligible embedded surface
field, including fields whose candidate is zero.  At either end
of an ordinary four-dimensional link with common base of dimension
at most one, its value is zero.
\end{theorem}

\begin{proof}
If there is a nonzero candidate, Proposition~\ref{sf:unique} says
that its eligible field is the only eligible field.  If there is
no nonzero candidate, all eligible candidates are zero and agree
with the definition.  In particular a zero candidate cannot be
ignored in favour of a nonzero candidate from another field.
The last assertion follows from Proposition~\ref{sf:conic-end},
Lemma~\ref{sf:surface-end}, and the definition when $\dim B\leq1$.
\end{proof}

\begin{theorem}[Vanishing between Mori spaces over points]
\label{sf:vanishing}
There is a bound $d_0$ with the following property.  Let $X_1$ and
$X_2$ be normal projective $\Q$-factorial terminal rationally connected
complex fourfolds whose maps to $\operatorname{Spec}\C$ are Mori
fibre spaces.  For $d>d_0$, every birational map
$f:X_1\dashrightarrow X_2$ satisfies
\[
                              c_u(f)=0.
\]
The bound is independent of $X_1$, $X_2$, $f$, its Sarkisov
factorization, and the chosen degree-$d$ test subcollection.
\end{theorem}

\begin{proof}
Choose an ordinary Sarkisov factorization supplied by
Theorem~\ref{bd:diagrams}, and orient each link from
$V/B$ to $V'/B'$.  Write $\Delta$ for final minus initial value of
a function on its nodes.  Proposition~\ref{sl:vertical} gives
\begin{equation}\label{sf:residual}
 c_u(\text{link})-\Delta v
 =H_Q=c_{u,\mathrm{hor}/Q}-\Delta h_Q.
\end{equation}
We check $H_Q=\Delta s$ in every dimension of the common base $Q$.

If $\dim Q\leq1$, Proposition~\ref{sl:low-base} gives $H_Q=0$, while
Theorem~\ref{sf:intrinsic} gives $s=0$ at both ends.

If $\dim Q=3$, both endpoint bases have dimension three and are
birational to $Q$.  On the generic fibre the link is an isomorphism
of smooth projective conics.  Thus the endpoint base fields and
quaternion classes are identified, including all eligible embedded
surface fields and their candidates.  Therefore $\Delta s=0$,
which equals $H_Q$ by Proposition~\ref{sl:low-base}.

Suppose $\dim Q=2$ and put $k=\C(Q)$.  In case~\textup{(ii)} of
Lemma~\ref{sl:generic-dichotomy}, the calculation of
Proposition~\ref{sl:surface-identity} gives
$c_{u,\mathrm{hor}/Q}=\Delta w$.  Here $w$ is evaluated on the
structured generic surface over $k$, including its conic structure
when $\dim B=3$, as in Definition~\ref{sl:w-def}.
On a del Pezzo side,
$h_Q=0$ and $s=w$.  On a conic side, its generic curve base over
$k$ is geometrically integral of genus zero.  Lemma~\ref{sl:conic-h}
identifies its singular-fibre set and component cover with
$\Delta_k$ and the residue cover.  The common del Pezzo surface $W$
has positive degree, so \eqref{sl:conic-size} gives
$r_k=\deg_k\Delta_k=8-K_W^2\leq7$.  Hence $k$ is eligible.
The same lemma gives $h_Q=-u(\Delta_k)$, and by
\eqref{sf:candidate} and Theorem~\ref{sf:intrinsic},
\[
                         s=w-h_Q.
\]
This remains valid if the candidate is zero.  Consequently
$H_Q=\Delta(w-h_Q)=\Delta s$.

In case~\textup{(i)} of Lemma~\ref{sl:generic-dichotomy}, the
structured generic surfaces are isomorphic after cancelling any
identical generic contractions.  Their exceptional horizontal
contribution and $H_Q$ are zero by that lemma.  At a
del Pezzo node the Picard-class sets defining $w$ identify.  At a
conic node the isomorphism identifies the embedded base field
and its quaternion class, so the intrinsic corrections identify.
Thus $\Delta s=0$ here as well.  This last argument requires no
bound on the discriminant degree of an isomorphic conic
structure.

It follows from \eqref{sf:residual} that every link satisfies
\[
                       c_u(\text{link})=\Delta(v+s).
\]
Add over the factorization and use the cocycle property
of Lemma~\ref{pre:cocycle}.  At a node over a point there
are no vertical divisors in the defining difference for $v$,
and $s=0$ by definition.  The sum therefore vanishes.  All lower
bounds on $d$ arose from the finitely many fixed families specified
after \eqref{sf:c-choice}; none depends on the length or starting
resolution of the factorization.
\end{proof}
